\documentclass[12pt]{article}
\usepackage[utf8]{inputenc}
\usepackage[T1]{fontenc}
\usepackage{hyperref}
\title{Sadri Language Resources for K-12 Curriculum Development}
\author{Compiled by Claude Sonnet 5.5 (claude-sonnet-5-5) for Andr\'as Kornai}
\date{}
\begin{document}
\maketitle

\section{Introduction}
Sadri (ISO 639-3 sck; Glottocode sadr1248), also called Nagpuri, Sadani or Chotanagpuri Sadri, is an Indo-Aryan lingua franca of eastern central India (Jharkhand and neighboring states). In the literary tradition it is called Nagpuri, and Nagpuri is the official name in Jharkhand. Speakers distinguish registers: Sadani (formal), common Sadri (Nagpuri), and lower Sadri. It is written mainly in Devanagari; Kaithi was used historically.

\section{Primary Educational Resources}
Searches found no Sadri primers, textbooks or school programmes. These would have to be sought through Jharkhand education bodies, which I did not find documented in results.

\section{Linguistic and Reference Materials}
\begin{itemize}
\item Peterson and Baraik, \emph{A Grammar of Chotanagpuri Sadri} (book listing), a recent grammar based largely on recordings of spontaneous speech. \url{https://www.exoticindiaart.com/book/details/grammar-of-chotanagpuri-sadri-indo-aryan-lingua-franca-of-eastern-central-india-hbh466}
\item Preliminary Sadri dictionary on Webonary, explicitly a work in progress. \url{https://www.webonary.work/sadri/?p=9444}
\item Indology list post (2019) on older short Sadani grammars (Whitley 1896, Floor 1931). \url{https://list.indology.info/pipermail/indology/2019-November/051286.html}
\item Wikipedia, ``Nagpuri language'': overview, names and scripts. \url{https://en.wikipedia.org/wiki/Nagpuri_language}
\item Simple English Wikipedia, ``Sadri language''. \url{https://simple.wikipedia.com/wiki/Sadri_language}
\item Omniglot page on Sadri writing. \url{https://omniglot.com/writing/sadri.htm}
\end{itemize}

\section{Religious and Community Resources}
\begin{itemize}
\item Scripture Earth page for Sadri: New Testament text and audio, with Gospel films. \url{https://www.scriptureearth.org/00eng.php?iso=sck}
\end{itemize}

\section{Digital and Audio Resources}
\begin{itemize}
\item Wikitongues page for Sadri (sck). \url{https://wikitongues.org/languages/sck/}
\item Scripture audio via the Scripture Earth page above.
\end{itemize}

\section{Curriculum Development Context}
The language is not known to be taught in schools, and I found no textbooks. Usable foundations are the Peterson and Baraik grammar, the preliminary Webonary dictionary and the scripture texts. A curriculum would need a decision on the standard name and register (Sadri versus Nagpuri) and on script (Devanagari), and would need coordination with Jharkhand authorities. Speaker numbers (about 3.5 million) are an estimate.
\end{document}